\documentclass[a4paper,12pt]{article}
\usepackage[ngerman]{babel}
\usepackage[utf8]{inputenc}
\usepackage{amsmath,amssymb}
\usepackage{tikz}
\usetikzlibrary{automata,positioning,arrows.meta}

\tikzset{
    ->,>=Stealth,
    node distance=2.5cm,
    every state/.style={thick,minimum size=1cm},
    initial text={},
}

\title{ETI -- Aufgabenblatt 1 \\ Hausaufgaben}
\author{}
\date{}

\begin{document}
\raggedright
\setlength{\parindent}{0pt}

{\LARGE\bfseries ETI -- Aufgabenblatt 1 \\ Hausaufgaben\par}
\vspace{2em}

\section*{Aufgabe 1.5}

\textbf{Sprache:} $L_1 = \{w \in \{a,b,c\}^* \mid \text{kein } c \text{ steht vor einem } a\}$.

Das bedeutet: Sobald ein $c$ gelesen wurde, darf danach kein $a$ mehr kommen. Die Wörter bestehen also aus einem $\{a,b\}$-Präfix gefolgt von einem $\{b,c\}$-Suffix.

\textbf{DFA:}

\begin{flushleft}
\begin{tikzpicture}
    \node[state, initial, accepting] (q0) {$q_0$};
    \node[state, accepting, right=of q0] (q1) {$q_1$};
    \node[state, right=of q1] (q2) {$q_2$};
    \path (q0) edge[loop above] node {$a,b$} ()
          (q0) edge node[above] {$c$} (q1)
          (q1) edge[loop above] node {$b,c$} ()
          (q1) edge node[above] {$a$} (q2)
          (q2) edge[loop above] node {$a,b,c$} ();
\end{tikzpicture}
\end{flushleft}

\textbf{Korrektheit:}
\begin{itemize}
    \item $q_0$: Startzustand -- es wurde noch kein $c$ gelesen. Alle $a$'s und $b$'s sind erlaubt.
    \item $q_1$: Es wurde mindestens ein $c$ gelesen, aber noch kein $a$ danach. Weitere $b$'s und $c$'s sind erlaubt.
    \item $q_2$: Falle -- es wurde ein $a$ nach einem $c$ gelesen. Das Wort wird abgelehnt.
\end{itemize}
$q_0$ und $q_1$ sind akzeptierend, $q_2$ nicht. Der Automat akzeptiert genau die Wörter, in denen kein $c$ vor einem $a$ steht.

\section*{Aufgabe 1.6}

\textbf{Sprache:} $L_2 = \{w \in \{a,\ldots,z\}^* \mid w \text{ enthält den Teilstring } eti\}$.

\subsection*{1. DFA $M$:}

Zustände: $q_0$ (Start), $q_1$ (``e'' gesehen), $q_2$ (``et'' gesehen), $q_3$ (``eti'' gesehen, akzeptierend).

\begin{flushleft}
\begin{tikzpicture}[node distance=3cm, every edge/.append style={font=\small},
    lbl/.style={fill=white, inner sep=2pt}]
    \node[state, initial left] (q0) {$q_0$};
    \node[state, below right=1.5cm and 3cm of q0] (q1) {$q_1$};
    \node[state, above right=1.5cm and 3cm of q1] (q2) {$q_2$};
    \node[state, accepting, below right=1.5cm and 3cm of q2] (q3) {$q_3$};
    % Forward edges
    \path (q0) edge node[lbl] {$e$} (q1)
          (q1) edge node[lbl] {$t$} (q2)
          (q2) edge node[lbl] {$i$} (q3);
    % Self-loops
    \path (q0) edge[loop above] node {$\Sigma \setminus \{e\}$} ()
          (q1) edge[loop below] node {$e$} ()
          (q3) edge[loop below] node {$\Sigma$} ();
    % Back edges
    \path (q1) edge[bend right=20] node[lbl] {$\Sigma \setminus \{e,t\}$} (q0);
    \path (q2) edge[bend right=20] node[lbl] {$e$} (q1);
    \path (q2) edge node[lbl] {$\Sigma \setminus \{e,i\}$} (q0);
\end{tikzpicture}
\end{flushleft}

\subsection*{2. Beweis $L(M) = L_2$:}

\textbf{Akzeptanz:}
Sei $M = (Q,\Sigma,\delta,q_0,F)$ mit $\Sigma = \{a,\ldots,z\}$ und $F = \{q_3\}$.
Sei außerdem $w = w_1 w_2 \cdots w_n$ eine Zeichenkette über $\Sigma$.
$M$ akzeptiert $w$, falls eine Folge von Zuständen $r_0,r_1,\ldots,r_n \in Q$ existiert mit:
\begin{itemize}
    \item $r_0 = q_0$
    \item $r_{i+1} = \delta(r_i,w_{i+1})$ für $i = 0,\ldots,n-1$
    \item $r_n \in F$
\end{itemize}

\textbf{Erkannte Sprache:}
\[
L(M) = \{w \in \Sigma^* \mid M \text{ akzeptiert } w\}.
\]

\textbf{Zu zeigen:} $L(M) = L_2$.

\medskip
\textbf{``$L(M) \subseteq L_2$'':}
Sei $w \in L(M)$. Dann akzeptiert $M$ das Wort $w$.
Nach Definition gibt es also eine Zustandsfolge $r_0,r_1,\ldots,r_n$ mit $r_0=q_0$, $r_{i+1}=\delta(r_i,w_{i+1})$ und $r_n \in F$.
Da $F=\{q_3\}$ gilt $r_n=q_3$.

Der Zustand $q_3$ wird nur durch den Übergang
\[
q_0 \xrightarrow{e} q_1 \xrightarrow{t} q_2 \xrightarrow{i} q_3
\]
erreicht.
Also muss in $w$ die Zeichenfolge ``eti'' vorkommen.
Damit gilt $w \in L_2$ und somit $L(M) \subseteq L_2$.

\medskip
\textbf{``$L_2 \subseteq L(M)$'':}
Sei $w \in L_2$. Dann enthält $w$ den Teilstring ``eti''.
Also lässt sich $w$ schreiben als
\[
w = x \cdot e \cdot t \cdot i \cdot y
\]
für $x,y \in \Sigma^*$.

Nach dem Lesen von $x$ befindet sich $M$ in einem Zustand aus $Q$.
Falls dieser Zustand bereits $q_3$ ist, wird $w$ akzeptiert, da $q_3$ für jedes Zeichen aus $\Sigma$ eine Schleife hat.
Andernfalls ist der Zustand $q_0$, $q_1$ oder $q_2$.
Beim Lesen des nächsten Zeichens $e$ geht $M$ in jedem dieser Fälle nach $q_1$.
Danach gilt:
\[
q_1 \xrightarrow{t} q_2 \xrightarrow{i} q_3.
\]
Der Rest $y$ wird in $q_3$ gelesen, und wegen der $\Sigma$-Schleife bleibt $M$ in $q_3$.
Also endet die Zustandsfolge in einem akzeptierenden Zustand.
Damit akzeptiert $M$ das Wort $w$, also $w \in L(M)$.
Somit gilt $L_2 \subseteq L(M)$.

\medskip
Aus $L(M) \subseteq L_2$ und $L_2 \subseteq L(M)$ folgt $L(M)=L_2$.

\section*{Aufgabe 1.7}

\textbf{Sprache:} $L_3 = \{s \in \{0,1\}^* \mid \sum_{i=1}^{|s|} s_i \cdot i \text{ ist gerade}\}$.

Wir verfolgen die Parität der Summe. Sei $S_k = \sum_{i=1}^{k} s_i \cdot i$ der Wert nach dem Lesen von $k$ Zeichen. Beim Lesen des $(k+1)$-ten Zeichens $s_{k+1}$ ändert sich die Parität von $S$ genau dann, wenn $s_{k+1} \cdot (k+1)$ ungerade ist, also wenn $s_{k+1} = 1$ \emph{und} $k+1$ ungerade.

Die Position $k+1$ ist ungerade für $k$ gerade und gerade für $k$ ungerade. Also brauchen wir als Zustand die Parität der Summe und die Parität der aktuellen Position.

\textbf{Zustände:} $(p, g)$ wobei $p \in \{0,1\}$ die Parität der Summe ist ($0$ = gerade) und $g \in \{0,1\}$ die Parität der aktuellen Position (nächstes Zeichen hat Position $g+1$, also: $g=0$ bedeutet nächste Position ist ungerade).

\textbf{Übergänge:} Beim Lesen von $s_{k+1}$ in Zustand $(p,g)$:
\begin{itemize}
    \item Lesen von $0$: Summe ändert sich nicht $\to (p,\, 1-g)$
    \item Lesen von $1$: Summe ändert sich um $k+1$. Falls $g=0$ (Position ungerade): Parität kippt $\to (1-p,\, 1)$. Falls $g=1$ (Position gerade): Parität bleibt $\to (p,\, 0)$.
\end{itemize}

\textbf{DFA:}

\begin{flushleft}
\begin{tikzpicture}[node distance=3cm]
    \node[state, initial, accepting] (00) {$(0,0)$};
    \node[state, accepting, right=of 00] (01) {$(0,1)$};
    \node[state, below=2.5cm of 00] (10) {$(1,0)$};
    \node[state, right=of 10] (11) {$(1,1)$};
    % Lesen von 0: g kippt, p bleibt
    \path (00) edge[bend left=15] node[above] {$0$} (01)
          (01) edge[bend left=15] node[below] {$0$} (00)
          (10) edge[bend left=15] node[above] {$0$} (11)
          (11) edge[bend left=15] node[below] {$0$} (10);
    % Lesen von 1: hängt von g ab
    % g=0: p kippt, g -> 1
    \path (00) edge node[right] {$1$} (11)
          (10) edge node[left] {$1$} (01);
    % g=1: p bleibt, g -> 0
    \path (01) edge[loop right] node {$1$} ()
          (11) edge[loop right] node {$1$} ();
\end{tikzpicture}
\end{flushleft}

\textbf{Formale Beschreibung:}
\[
M = (Q,\Sigma,\delta,(0,0),F)
\]
mit
\[
Q=\{(0,0),(0,1),(1,0),(1,1)\}, \quad \Sigma=\{0,1\}, \quad F=\{(0,0),(0,1)\}.
\]

\textbf{Akzeptanz:}
Sei $s = s_1s_2\cdots s_n$ eine Zeichenkette über $\Sigma$.
$M$ akzeptiert $s$, falls eine Folge von Zuständen $r_0,r_1,\ldots,r_n \in Q$ existiert mit:
\begin{itemize}
    \item $r_0 = (0,0)$
    \item $r_{i+1} = \delta(r_i,s_{i+1})$ für $i = 0,\ldots,n-1$
    \item $r_n \in F$
\end{itemize}

\textbf{Erkannte Sprache:}
\[
L(M)=\{s \in \Sigma^* \mid M \text{ akzeptiert } s\}.
\]

\textbf{Zu zeigen:} $L(M)=L_3$.

\medskip
\textbf{Korrektheit:}
Wir zeigen per Induktion über $k$, dass nach dem Lesen der ersten $k$ Zeichen gilt:
\[
r_k=(p,g) \quad \Longleftrightarrow \quad
p \equiv \sum_{i=1}^{k} s_i \cdot i \pmod 2
\text{ und } g \equiv k \pmod 2.
\]

\textbf{Induktionsanfang:}
Für $k=0$ ist $r_0=(0,0)$.
Die leere Summe ist $0$ und es wurden $0$ Zeichen gelesen.
Also gilt die Behauptung.

\textbf{Induktionsschritt:}
Sei die Behauptung für $k$ richtig und sei $r_k=(p,g)$.
Wir betrachten das nächste Zeichen $s_{k+1}$.
\begin{itemize}
    \item Falls $s_{k+1}=0$, ändert sich die Summe nicht. Nur die Parität der gelesenen Zeichen wechselt. Also geht $M$ nach $(p,1-g)$.
    \item Falls $s_{k+1}=1$ und $g=0$, dann ist $k$ gerade und $k+1$ ungerade. Die Summe ändert ihre Parität, also geht $M$ nach $(1-p,1)$.
    \item Falls $s_{k+1}=1$ und $g=1$, dann ist $k$ ungerade und $k+1$ gerade. Die Summe ändert ihre Parität nicht, also geht $M$ nach $(p,0)$.
\end{itemize}
Damit speichert der Zustand nach jedem gelesenen Präfix genau die Parität der Summe und die Parität der Anzahl gelesener Zeichen.

Am Ende gilt also für $r_n=(p,g)$:
\[
p \equiv \sum_{i=1}^{n} s_i \cdot i \pmod 2.
\]
Der Automat akzeptiert genau dann, wenn $r_n \in F$ gilt.
Da $F=\{(0,0),(0,1)\}$, ist das genau dann der Fall, wenn $p=0$ ist.
Also akzeptiert $M$ genau die Wörter, für die
\[
\sum_{i=1}^{|s|} s_i \cdot i
\]
gerade ist.
Damit gilt $L(M)=L_3$.

\end{document}
